\chapter{Równanie Schrödingera w trzech wymiarach}

\begin{summarybox}
Jednowymiarowe modele uczą podstaw, ale rzeczywiste układy kwantowe są zwykle
trójwymiarowe. W trzech wymiarach pojawia się laplasjan, geometria sferyczna,
moment pędu oraz radialna gęstość prawdopodobieństwa. Ten rozdział pokazuje, jak
przejść od równania 1D do równania Schrödingera w 3D i jak rozumieć separację
zmiennych w potencjałach centralnych.
\end{summarybox}

\section{Od problemu 1D do problemu 3D}

W jednym wymiarze stacjonarne równanie Schrödingera ma postać
\begin{equation}
    -\frac{\hbar^2}{2m}\frac{d^2\psi}{dx^2}+V(x)\psi=E\psi.
\end{equation}
W trzech wymiarach funkcja falowa zależy od trzech współrzędnych:
\begin{equation}
    \psi=\psi(x,y,z),
\end{equation}
a druga pochodna po \(x\) zostaje zastąpiona przez laplasjan:
\begin{equation}
    \nabla^2=\frac{\partial^2}{\partial x^2}+\frac{\partial^2}{\partial y^2}+\frac{\partial^2}{\partial z^2}.
\end{equation}
Równanie stacjonarne brzmi więc
\begin{equation}
    -\frac{\hbar^2}{2m}\nabla^2\psi(\mathbf{r})+V(\mathbf{r})\psi(\mathbf{r})=E\psi(\mathbf{r}).
\end{equation}

Interpretacja probabilistyczna też się zmienia. Prawdopodobieństwo znalezienia cząstki w objętości \(dV\) wynosi
\begin{equation}
    |\psi(\mathbf{r})|^2\,dV.
\end{equation}
Warunek normalizacji ma postać
\begin{equation}
    \int |\psi(\mathbf{r})|^2\,d^3r=1.
\end{equation}

\section{Współrzędne kartezjańskie i sferyczne}

Współrzędne kartezjańskie \((x,y,z)\) są naturalne dla potencjałów o prostych ścianach, pudełek i siatek numerycznych. Współrzędne sferyczne \((r,\theta,\phi)\) są naturalne dla potencjałów centralnych, czyli takich, które zależą tylko od odległości od początku układu:
\begin{equation}
    V(\mathbf{r})=V(r).
\end{equation}

Związek między współrzędnymi jest następujący:
\begin{align}
    x&=r\sin\theta\cos\phi,\\
    y&=r\sin\theta\sin\phi,\\
    z&=r\cos\theta.
\end{align}
Element objętości we współrzędnych sferycznych wynosi
\begin{equation}
    d^3r=r^2\sin\theta\,dr\,d\theta\,d\phi.
\end{equation}
Ten czynnik \(r^2\sin\theta\) jest bardzo ważny przy normalizacji i interpretacji radialnej gęstości prawdopodobieństwa.

\section{Laplasjan we współrzędnych sferycznych}

Laplasjan w układzie sferycznym ma postać
\begin{equation}
    \nabla^2=\frac{1}{r^2}\frac{\partial}{\partial r}\left(r^2\frac{\partial}{\partial r}\right)
    +\frac{1}{r^2\sin\theta}\frac{\partial}{\partial\theta}\left(\sin\theta\frac{\partial}{\partial\theta}\right)
    +\frac{1}{r^2\sin^2\theta}\frac{\partial^2}{\partial\phi^2}.
\end{equation}
Część kątową można zapisać za pomocą operatora momentu pędu:
\begin{equation}
    \hat{L}^2=-\hbar^2\left[
    \frac{1}{\sin\theta}\frac{\partial}{\partial\theta}\left(\sin\theta\frac{\partial}{\partial\theta}\right)
    +\frac{1}{\sin^2\theta}\frac{\partial^2}{\partial\phi^2}
    \right].
\end{equation}
Dlatego Hamiltonian dla potencjału centralnego można zapisać jako
\begin{equation}
    \hat{H}=-\frac{\hbar^2}{2m}\frac{1}{r^2}\frac{\partial}{\partial r}\left(r^2\frac{\partial}{\partial r}\right)
    +\frac{\hat{L}^2}{2mr^2}+V(r).
\end{equation}

\section{Separacja zmiennych}

Dla potencjału centralnego szukamy rozwiązań w postaci
\begin{equation}
    \psi(r,\theta,\phi)=R(r)Y_l^m(\theta,\phi).
\end{equation}
Część kątowa jest harmoniką sferyczną, ponieważ spełnia równanie własne
\begin{equation}
    \hat{L}^2Y_l^m=\hbar^2l(l+1)Y_l^m.
\end{equation}
Wtedy problem trójwymiarowy redukuje się do równania radialnego dla \(R(r)\).

Separacja zmiennych działa dzięki symetrii. Jeśli potencjał zależy tylko od \(r\), to kierunki w przestrzeni są równoważne, a moment pędu jest dobrym narzędziem klasyfikacji stanów.

\section{Część radialna i część kątowa}

Po podstawieniu \(\psi=RY\) dostajemy równanie radialne. Wygodnie wprowadzić funkcję
\begin{equation}
    u(r)=rR(r).
\end{equation}
Wtedy równanie radialne przyjmuje postać podobną do jednowymiarowego równania Schrödingera:
\begin{equation}
    -\frac{\hbar^2}{2m}\frac{d^2u}{dr^2}
    +\left[V(r)+\frac{\hbar^2l(l+1)}{2mr^2}\right]u=Eu.
\end{equation}
Drugi składnik w nawiasie to bariera odśrodkowa. Dla \(l>0\) utrudnia ona znalezienie cząstki blisko środka układu.

Część kątowa opisuje rozkład kierunkowy, a część radialna opisuje zależność od odległości. Pełna gęstość prawdopodobieństwa wynosi
\begin{equation}
    |\psi(r,\theta,\phi)|^2=|R(r)|^2|Y_l^m(\theta,\phi)|^2.
\end{equation}

\section{Warunki regularności w początku układu}

Funkcja falowa musi być skończona i fizycznie dopuszczalna w punkcie \(r=0\). Dla funkcji \(u(r)=rR(r)\) warunek regularności zwykle oznacza
\begin{equation}
    u(0)=0.
\end{equation}
Jeśli \(u(0)\neq0\), to \(R(r)=u(r)/r\) byłoby osobliwe jak \(1/r\). W wielu problemach taka osobliwość jest niedopuszczalna.

Dla małych \(r\), rozwiązania radialne zachowują się typowo jak
\begin{equation}
    R(r)\sim r^l.
\end{equation}
Oznacza to, że dla większego \(l\) funkcja falowa jest silniej tłumiona w pobliżu środka.

\section{Warunki brzegowe na skończonym promieniu}

Jeżeli cząstka jest zamknięta w sferycznej studni o promieniu \(R\), często stosujemy warunek
\begin{equation}
    \psi(R,\theta,\phi)=0.
\end{equation}
Dla części radialnej oznacza to
\begin{equation}
    R(R)=0
\end{equation}
albo równoważnie
\begin{equation}
    u(R)=0.
\end{equation}
Warunek ten prowadzi do kwantowania energii, podobnie jak warunki brzegowe w jednowymiarowej studni.

Dla potencjałów zanikających w nieskończoności, jak w atomie wodoru, zamiast warunku na skończonym promieniu wymagamy normalizowalności:
\begin{equation}
    \int_0^\infty |R(r)|^2r^2\,dr<\infty.
\end{equation}

\section{Interpretacja radialnej gęstości prawdopodobieństwa}

Gęstość \(|\psi|^2\) jest gęstością na jednostkę objętości. Jeśli pytamy o prawdopodobieństwo znalezienia cząstki między \(r\) i \(r+dr\), musimy uwzględnić objętość powłoki sferycznej. Dla znormalizowanych harmonik sferycznych radialna gęstość prawdopodobieństwa wynosi
\begin{equation}
    P(r)=r^2|R(r)|^2.
\end{equation}
W ujęciu przez \(u(r)=rR(r)\) mamy po prostu
\begin{equation}
    P(r)=|u(r)|^2.
\end{equation}

To wyjaśnia częsty błąd interpretacyjny: maksimum \(|R(r)|^2\) nie musi oznaczać najbardziej prawdopodobnego promienia. Najbardziej prawdopodobny promień wyznacza maksimum \(r^2|R(r)|^2\).

\begin{summarybox}
W trzech wymiarach funkcja falowa zależy od geometrii przestrzeni. Dla potencjałów centralnych naturalny jest rozkład na część radialną i kątową. Moment pędu opisuje część kątową, a radialna gęstość prawdopodobieństwa wymaga czynnika \(r^2\).
\end{summarybox}
